\subsection{HALL BASES OF A FREE LIE ALGEBRA}

We preserve the notation of the preceding no.

THEOREM 1. Let H be a Hall set relative to X and \(\Psi\) the canonical mapping of M(X) into the free Lie algebra L(X). The restriction of \(\Psi\) to H is a basis of the module L(X).

For every element \(w\) of H we write \(\bar{w} = \Psi(w)\) .

\begin{enumerate}
\def\labelenumi{(\Alph{enumi})}
\tightlist
\item
  The case where X is finite.
\end{enumerate}

If X is empty, so are \(M(X)\) and therefore H and \(L(X)\) is zero. If X has a single element x, \(H \cap M^{n}(X)\) is empty for \(n \geqslant 2\) (Proposition 12 (a)). Therefore, \(H = \{x\}\) ; we know also (no. 2, Remark) that the module \(L(X)\) is free and has basis \(\{\bar{x}\}\) . The theorem is therefore true when X has at most one element.

Suppose henceforth that X has at least two elements; choose sequences \((w_{p})\) and \((P_{p})\) with the properties stated in Proposition 13. For every integer \(p \geqslant 0\) , let \(L_{p}\) denote the submodule of \(\mathrm{L}(\mathbf{X})\) generated by the elements \(\bar{w}_{i}\) with \(0 \leqslant i < p\) and \(g_{p}\) the Lie subalgebra of \(\mathrm{L}(\mathbf{X})\) generated by the family \((\bar{u})_{u \in \mathbf{P}_{p}}\) .

Lemma 2. For every integer \(p \geqslant 0\) , the module \(\mathrm{L}_p\) admits the family \((\bar{w}_i)_{0 \leqslant i < p}\) as basis, the Lie algebra \(\mathfrak{g}_p\) admits \((\bar{u})_{u \in \mathbf{P}_p}\) as basic family and the module \(\mathrm{L}(\mathbf{X})\) is the direct sum of \(\mathrm{L}_p\) and \(\mathfrak{g}_p\) .

\(L_{0} = \{0\}\) and \(g_{0} = L(X)\) and the lemma is true for \(p = 0\) . We argue by induction on \(p\) . Suppose then that the lemma is true for some integer \(p \geqslant 0\) . Let \(u_{i,w} = (\text{ad } \bar{w}_{p})^{i}. \bar{w} = \Psi(w_{p}^{i} w)\) for \(i \geqslant 0, w \in P_{p}, w \neq w_{p}\) . By the Corollary to Proposition 10 of no. 9, the free Lie algebra \(g_{p}\) is the direct sum of the module \(T_{p}\) of basis \(\{\bar{w}_{p}\}\) and a Lie subalgebra \(b_{p}\) admitting

\[
\mathcal {F} = (u _ {i, w}) _ {i \geqslant 0, w \in \mathrm{P} _ {p}, w \neq w _ {p}}
\]

as basic family. By Proposition 13 (d), the family \((\bar{u})_{u\in \mathbf{P}_{p + 1}}\) is equal to \(\mathcal{F}\) and hence is a basic family of \(\mathfrak{h}_p = \mathfrak{g}_{p + 1}\) . Hence \(\mathrm{L}(\mathrm{X}) = \mathrm{L}_p\oplus \mathrm{T}_p\oplus \mathfrak{g}_{p + 1}\) and, as \(\mathrm{L}_{p + 1} = \mathrm{L}_p + \mathrm{T}_p\) , \(\mathrm{L}(\mathrm{X}) = \mathrm{L}_{p + 1}\oplus \mathfrak{g}_{p + 1}\) and \((\bar{w}_0,\bar{w}_1,\dots ,\bar{w}_{p - 1},\bar{w}_p)\) is a basis of the module \(\mathbf{L}_{p + 1}\) .

Let \(n\) be a positive integer. By Proposition 13 (c) there exists an integer \(p(n)\) such that \(P_p\) has only elements of length \(>n\) for \(p \geqslant p(n)\) . For \(p \geqslant p(n)\) , the Lie subalgebra \(g_p\) of \(L(X)\) is generated by elements of degree \(>n\) and hence \(L^n(X) \cap g_p = \{0\}\) . On the other hand, the elements \(\bar{w}_i\) of \(L(X)\) are homogeneous and the family \((\bar{w}_i)_{0 \leqslant i < p}\) is a basis of a supplementary module of \(g_p\) . It follows immediately that the family of elements \(\bar{w}_i\) of degree \(n\) is a basis of the module \(L^n(X)\) and that the sequence \((\bar{w}_i)_{i \geqslant 0}\) is a basis of the module \(L(X)\) .

\begin{enumerate}
\def\labelenumi{(\Alph{enumi})}
\setcounter{enumi}{1}
\tightlist
\item
  General case.
\end{enumerate}

If \(S\) is a subset of \(X\) , recall that \(M(S)\) is identified with the submagma of \(M(X)\) generated by \(S\) and \(L(S)\) is identified with the Lie subalgebra of \(L(X)\) generated by \(S\) ; we have seen that if \(w \in M(S)\) is of length \(\geqslant 2\) then \(\alpha(w) \in M(S)\) and \(\beta(w) \in M(S)\) . It follows immediately that \(H \cap M(S)\) is a Hall set relative to \(S\) .

For every finite subset \(\Phi\) of \(H\) there exists a finite subset \(S\) of \(X\) such that \(\Phi \subset \mathrm{M}(\mathrm{S})\) . Case (A) then shows that the elements \(\bar{w}\) with \(w \in \Phi\) are linearly independent in \(\mathrm{L}(\mathrm{S})\) and hence in \(\mathrm{L}(\mathrm{X})\) . Therefore the family \((\bar{w})_{w \in \mathbf{H}}\) is free.

For every element \(a\) of \(\mathrm{L}(\mathrm{X})\) there exists a finite subset \(\mathrm{S}\) of \(\mathrm{X}\) such that \(a \in \mathrm{L}(\mathrm{S})\) . By case (A), the subset \(\Psi(\mathrm{H} \cap \mathrm{M}(\mathrm{S}))\) of \(\Psi(\mathrm{H})\) generates the module \(\mathrm{L}(\mathrm{S})\) and hence \(a\) is a linear combination of elements of \(\Psi(\mathrm{H})\) . Hence \(\Psi(\mathrm{H})\) generates the module \(\mathrm{L}(\mathrm{X})\) , which completes the proof.

COROLLARY. The module \(\mathbf{L}(\mathbf{X})\) is free and so is each of the submodules \(\mathbf{L}^{\alpha}(\mathbf{X})\) where \(\alpha \in \mathbf{N}^{(\mathbf{X})}\) and \(\mathbf{L}^n (\mathbf{X})\) where \(n\in \mathbf{N}\) . The modules \(\mathbf{L}^{\alpha}(\mathbf{X})\) are of finite rank and so are the modules \(\mathbf{L}^n (\mathbf{X})\) if \(\mathbf{X}\) is finite.

There exists a Hall set H relative to X (Proposition 11). For all \(w \in \mathbf{H}\) , the element \(\Psi(w)\) of \(\mathbf{L}(\mathbf{X})\) belongs to one of the modules \(\mathrm{L}^{\alpha}(\mathbf{X})\) (with \(\alpha \in \mathbf{N}^{(\mathbf{X})}\) ) and the module \(\mathrm{L}(\mathbf{X})\) is the direct sum of the submodules \(\mathrm{L}^{\alpha}(\mathbf{X})\) . Further, for all \(\alpha \in \mathbf{N}^{(\mathbf{X})}\) , the set of elements of \(\mathrm{M}(\mathbf{X})\) whose canonical image in \(\mathbf{N}^{(\mathbf{X})}\) is equal to \(\alpha\) is finite; this shows that each of the modules \(\mathrm{L}^{\alpha}(\mathbf{X})\) is free and of finite rank and that \(\mathrm{L}(\mathbf{X})\) is free. Now \(\mathrm{L}^{n}(\mathbf{X}) = \bigoplus_{|\alpha| = n} \mathrm{L}^{\alpha}(\mathbf{X})\) and hence \(\mathrm{L}^{n}(\mathbf{X})\) is free; when X is finite, the set of \(\alpha \in \mathbf{N}^{(\mathbf{X})}\) such that \(|\alpha| = n\) is finite and hence \(\mathrm{L}^{n}(\mathbf{X})\) is then of finite rank.

DEFINITION 3. A Hall basis of a free Lie algebra \(\mathbf{L}(\mathbf{X})\) is any basis of \(\mathbf{L}(\mathbf{X})\) which is the canonical image of a Hall set relative to \(\mathbf{X}\) .

Remark. Suppose that X consists of two distinct elements x and y and let \(L^{(\cdot,1)}\) be the submodule of \(L(X)\) the sum of the \(L^{\alpha}(X)\) where \(\alpha\in N^{X}\) and \(\alpha(y)=1\) . It follows immediately from Theorem 1 and Proposition 12 of no. 10 that the elements of \((ad x)^{n}\cdot y\) where n is an integer \(\geqslant0\) form a basis of the submodule \(L^{(\cdot,1)}\) . It follows that the restriction to \(L^{(\cdot,1)}\) of the mapping ad x is injective.

